% FONTE TEMA https://github.com/matze/mtheme
\documentclass[aspectratio=1610]{beamer}
%\documentclass[aspectratio=1610, handout]{beamer}
\usepackage[utf8]{inputenc}
\usepackage{ragged2e}
\usepackage{xcolor}
\usepackage[italian]{babel}
\usepackage{multirow}
\usepackage{silence}
\usepackage{listings}
\definecolor{verde}{rgb}{0,0.6,0}
\definecolor{grigio}{rgb}{0.5,0.5,0.5}
\definecolor{rosa}{rgb}{0.58,0,0.82}
\definecolor{backgroud}{rgb}{0.95,0.95,0.92}
\lstdefinestyle{stilepython}{
    backgroundcolor=\color{backgroud},
    commentstyle=\color{verde},
    keywordstyle=\color{magenta},
    numberstyle=\tiny\color{grigio},
    stringstyle=\color{rosa},
    basicstyle=\ttfamily\footnotesize,
    breakatwhitespace=false,
    numbers=left,
    numbersep=5pt,
    showspaces=false,
    showstringspaces=false,
    showtabs=false,
    literate=
    {à}{{\'a}}1
    {è}{{\'e}}1
    {ì}{{\'i}}1
    {ò}{{\'o}}1
    {ù}{{\'u}}1
}
\lstset{style=stilepython}

\WarningFilter{beamer}{}
\WarningFilter{metropolis}{}
\usetheme[progressbar=frametitle,titleformat=smallcaps]{metropolis}
\setbeamertemplate{frame numbering}[fraction]
\setbeamercovered{dynamic}
\definecolor{rosso}{RGB}{255, 0, 0}
\definecolor{giallo}{RGB}{254,212,23}
\hypersetup{colorlinks=true,linkcolor=black,urlcolor=rosso}
\setbeamercolor{palette primary}{fg=black, bg=giallo}
\setbeamercolor{background canvas}{bg=white}
\setbeamercolor{normal text}{fg=black}
\setbeamercolor{progress bar}{fg=rosso}
\setbeamercolor{framesubtitle}{fg=rosso}
\setbeamercolor{normal text .dimmed}{fg=giallo}
\setbeamercolor{block title alerted}{fg=rosso, bg=giallo}
\setbeamerfont{caption}{size=\tiny}
\setbeamerfont{caption name}{size=\tiny}
\setlength{\abovecaptionskip}{0pt}
\makeatletter
\metroset{block=fill}
\setlength{\metropolis@progressinheadfoot@linewidth}{1pt} 
\setlength{\metropolis@progressonsectionpage@linewidth}{1pt}
\setlength{\metropolis@titleseparator@linewidth}{1pt}
\makeatother

\title{INPUT E OUTPUT}
\subtitle{Funzioni Python per l'interazione utente}
\date{}
\institute{\textit{
        Fonti:
        \begin{itemize}
            \item[-] \href{https://docs.python.org/it/3/tutorial/inputoutput.html}{Documentazione Ufficiale Python}
            \item[-] \href{https://www.zanichelli.it/ricerca/prodotti/coding-python}{Coding con Python - Zanichelli}
        \end{itemize}
    }
}

\begin{document}

\begin{frame}[plain, noframenumbering]
    \titlepage
\end{frame}

\section{FUNZIONI DI INPUT}

\begin{frame}[fragile]{INPUT}
    \begin{alertblock}{CASTING DELL'INPUT}
        \begin{minipage}{0.98\linewidth}
            \justifying
            La funzione \textcolor{red}{\textbf{input(}}prompt\textcolor{red}{\textbf{)}} 
            sospende l'esecuzione del programma rimanendo in attesa dell'immissione di 
            dati da parte dell'utente; una volta confermato l'inserimento tramite il 
            tasto Invio, restituisce la riga acquisita come valore di tipo \textbf{str}, 
            mostrando preventivamente a schermo l'eventuale testo specificato nell'
            argomento opzionale prompt.\\
            \bigskip
\begin{lstlisting}[language=Python]
 nome = input("Inserisci il tuo nome: ")
\end{lstlisting}
        \end{minipage}
    \end{alertblock}
\end{frame}

\begin{frame}[fragile]{INPUT}
    \begin{alertblock}{CASTING DELL'INPUT}
        \begin{minipage}{0.98\linewidth}
\begin{lstlisting}[language=Python]
 # 1. Casting a Int (numero intero)
 eta = int(input("Inserisci la tua età: "))

 # 2. Casting a Float (Numero con la virgola)
 media = float(input("Inserisci la tua media voti (es. 5.75): "))
 
 # 3. Casting a Bool, ovvero, True nel caso sia presente un valore 
 # diverso da zero o diverso da null, False altrimenti
 risposta = bool(input("Premi invio per inserire NO, inserisci un 
 carattere qualsiasi per inserire SI: "))
\end{lstlisting}
        \end{minipage}
    \end{alertblock}
\end{frame}

\begin{frame}[fragile]{INPUT}
    \begin{alertblock}{FUNZIONE MAP() E METODO SPLIT()}
        \begin{minipage}{0.98\linewidth}
            \justifying
            Il \textbf{metodo .split()} separa la stringa dell'utente in più elementi dividendo per spazi, 
            mentre la \textbf{funzione map()} applica la conversione in intero (int) a ciascuno di essi; 
            usati insieme, permettono di estrarre, convertire e assegnare più numeri scritti su 
            un'unica riga alle variabili num1 e num2 in un solo passaggio.
            \bigskip
\begin{lstlisting}[language=Python]
 # Casting dell'input su più righe ripetute
 num1 = int(input("inserisci un numero: "))
 num2 = int(input("inserisci un numero: "))
 # Si può rendere più compatto utilizzando la funzione map() in
 # combinazione con input() e il metodo split() fornito dal tipo str
 num1, num2 = map(int, input("Inserisci due numeri: ").split())
\end{lstlisting}
        \end{minipage}
    \end{alertblock}
\end{frame}

\section{FUNZIONI DI OUTPUT}

\begin{frame}[fragile]{OUTPUT}
    \begin{alertblock}{FORMATTAZIONE DELL'OUTPUT}
        \begin{minipage}{0.98\linewidth}
            \justifying
            La funzione \textcolor{red}{\textbf{print(}}argomenti\textcolor{red}{\textbf{)}} è lo strumento 
            principale per inviare dati verso lo schermo (o un qualsiasi flusso di output). Il suo 
            funzionamento interno è lineare: prende un numero variabile di argomenti, li converte 
            in stringhe (richiamando implicitamente il metodo str()), li unisce inserendo un separatore 
            tra di essi, ed emette la stringa finale seguita da un carattere di chiusura.
            \bigskip
\begin{lstlisting}[language=Python]
 print(*objects, sep=' ', end='\n', file=sys.stdout, flush=False)
\end{lstlisting}
        \end{minipage}
    \end{alertblock}
\end{frame}

\begin{frame}[fragile]{OUTPUT}
    \begin{alertblock}{FORMATTAZIONE DELL'OUTPUT}
        \begin{minipage}{0.98\linewidth}
\begin{lstlisting}[language=Python]
 # 1. Stampa base e argomenti multipli
 nome, eta = "Lele", 32
 print("Ciao", nome, "hai", eta, "anni.")

 # 2. Personalizzazione dei parametri sep ed end
 print("1789", "07", "14", sep="-")

 # 3. Personalizzazione del parametro end
 print("Elaborazione in corso...", end=" ")
 print("Completata!")
\end{lstlisting}
        \end{minipage}
    \end{alertblock}
\end{frame}

\begin{frame}[fragile]{OUTPUT}
    \begin{alertblock}{FORMATTAZIONE DELL'OUTPUT}
        \begin{minipage}{0.98\linewidth}
\begin{lstlisting}[language=Python]
 # 3. Metodo f-Strings, confronto con stampa base e funzione format
 prod, p, sconto = "Laptop", 899.986, 0.15
 print(f"Prod: {prod} | Sconto: {p * (1 - sconto):.2f}")
 # Più comodo che utilizzare la seguente stampa base
 print("Prod:", prod, "| Sconto:", format(p * (1 - sconto), ".2f"))
\end{lstlisting}
        \end{minipage}
    \end{alertblock}
\end{frame}

\end{document}